%\documentclass{article}
\documentclass[letterpaper,10pt]{article}

% if you need to pass options to natbib, use, e.g.:
%\usePackage{numbers, compress}{natbib}
% before loading neurips_2023


% ready for submission
%\usepackage{neurips_2023}

% to compile a preprint version, e.g., for submission to arXiv, add add the
% [preprint] option:
\usepackage[preprint]{neurips_2023}
\usepackage{natbib}
\setcitestyle{numbers,square}
% to compile a camera-ready version, add the [final] option, e.g.:
%     \usepackage[final]{neurips_2023}
% to avoid loading the natbib package, add option nonatbib:
%\usepackage[nonatbib]{neurips_2023}
\usepackage{times}
\usepackage{latexsym}
\usepackage{xcolor}
\usepackage{makecell}
\usepackage{times}
\usepackage{hyperref}
\usepackage{inconsolata}
\usepackage{url}
\usepackage{multirow}
%\usepackage{graphicx}
\usepackage[pdftex]{graphicx}
\usepackage{caption}
\usepackage{subcaption}
\usepackage{amsmath}
\usepackage{bbm}
\usepackage{booktabs}
\usepackage{algorithmic}
\usepackage{amssymb}
\usepackage{bm}
\usepackage[ruled,vlined,linesnumbered]{algorithm2e}
\usepackage{enumitem}
\usepackage{pifont}
\usepackage{pbox}
\usepackage[utf8]{inputenc} % allow utf-8 input
\usepackage[T1]{fontenc}    % use 8-bit T1 fonts
%\usepackage{hyperref}       % hyperlinks
\usepackage{url}            % simple URL typesetting
\usepackage{booktabs}       % professional-quality tables
\usepackage{amsfonts}       % blackboard math symbols
\usepackage{nicefrac}       % compact symbols for 1/2, etc.
\usepackage{microtype}      % microtypography
\usepackage{xcolor}         % colors
% \usepackage{subfig}
\usepackage{tabularx}
\newcommand{\linjie}[1]{\textcolor{blue}{[~Linjie: #1~]}}
\newcommand{\kevin}[1]{\textcolor{cyan}{[~Kevin: #1~]}}
\newcommand{\fuxiao}[1]{\textcolor{red}{[~fuxiao: #1~]}}
\newcommand{\jianfw}[1]{\textcolor{pink}{[~Jianfeng: #1~]}}
\newcommand{\yy}[1]{\textcolor{green}{[~yaser: #1~]}}

\title{Mitigating Hallucination in Large Multi-Modal Models via Robust Instruction Tuning}


% \author{%
%   Fuxiao Liu\\
%   University of Maryland, College Park\\  
%   \texttt{fl3es@umd.edu} \\
%   Yaser Yacoob\\
%   University of Maryland, College Park\\  
%   \texttt{yaser@umd.edu} \\
%   % examples of more authors
%   % \And
%   % Coauthor \\
%   % Affiliation \\
%   % Address \\
%   % \texttt{email} \\
%   % \AND
%   % Coauthor \\
%   % Affiliation \\
%   % Address \\
%   % \texttt{email} \\
%   % \And
%   % Coauthor \\
%   % Affiliation \\
%   % Address \\
%   % \texttt{email} \\
%   % \And
%   % Coauthor \\
%   % Affiliation \\
%   % Address \\
%   % \texttt{email} \\
% }

\author{%
 Fuxiao Liu$^1$, Kevin Lin$^2$, Linjie Li$^2$, Jianfeng Wang$^2$, Yaser Yacoob$^1$, Lijuan Wang$^2$\\ 
$^1$University of Maryland, College Park\quad $^2$Microsoft Corporation\\
\texttt{\small\{fl3es, yaser\}@umd.edu}, 
\texttt{\small\{keli, lindsey.li, jianfw, lijuanw\}@microsoft.com}
}
\date{June 2023}
\begin{document}
\maketitle

\begin{abstract}
Despite the promising progress in multi-modal tasks, current large multi-modal models (LMM) are prone to hallucinating inconsistent descriptions with respect to the associated image and human instructions. 
%We are the first to build a large and diverse visual instruction tuning dataset to address this issue. Our dataset \textit{Large-scale Robust Visual (LRV)-Instruction} contains 120k visual instructions, with 16 vision and language tasks and negative instructions. 
This paper addresses this issue by introducing the first large and diverse visual instruction tuning dataset, named \textit{Large-scale Robust Visual (LRV)-Instruction}.
% The dataset consists of 120k visual instructions with 16 vision-and-language tasks, including both positive and negative instructions. 
Our dataset consists of 120k visual instructions generated by GPT4, covering 16 vision-and-language tasks with open-ended instructions and answers. 
Unlike existing studies that primarily focus on positive instruction samples, we design \textit{LRV-Instruction} to include both positive and negative instructions for more robust visual instruction tuning.
Our negative instructions are designed at two semantic levels: \textit{(i) Nonexistent Element Manipulation} and \textit{(ii) Existent Element Manipulation}. To efficiently measure the hallucination generated by LMMs, we propose \textit{GPT4-Assisted Visual Instruction Evaluation (GAVIE)}, a novel approach to evaluate visual instruction tuning without the need for human-annotated groundtruth answers and can adapt to diverse instruction formats. We conduct comprehensive experiments to investigate the hallucination of LMMs. 
%The experiments show that existing LMMs seriously hallucinate when presented with our negative instructions. They suffer more from \textit{Existent Element Manipulation} instructions than from \textit{Nonexistent Element Manipulation} instructions. 
Our results demonstrate that existing LMMs exhibit significant hallucination when presented with our negative instructions, particularly with Existent Element Manipulation instructions.
Moreover, by finetuning MiniGPT4 on \textit{LRV-Instruction}, 
%we not only largely mitigate the hallucination but also boost its performance on public datasets with much less training data than state-of-art methods. 
we successfully mitigate hallucination while improving performance on public datasets using less training data compared to state-of-the-art methods.
% We also found that a balanced ratio between positive and negative instances in training data will generate a more robust model.
Additionally, we observed that a balanced ratio of positive and negative instances in the training data leads to a more robust model. Updates of our project are available at this \href{https://fuxiaoliu.github.io/LRV/}{link}. 

\end{abstract}


\section{Introduction}
\begin{figure*}[t]
    \centering
      \includegraphics[width=1\textwidth]{model2.pdf}
    \caption{Given an image and human instruction as the input, we introduce \textit{GPT4-Assisted Visual Instruction Evaluation (GAVIE)} to assess the output from current LMMs, such as \textit{MiniGPT4 and LLaVA}. \textbf{\textcolor{blue}{BLUE}} represents LMMs can not accurately follow human instructions while \textbf{\textcolor{red}{RED}} means they suffer from the hallucination problem. We can generate a more robust answer after finetuning MiniGPT4 on our proposed \textit{LRV-Instruction} dataset.  }
    \label{fig:model}
    \vspace{-0.15in}
\end{figure*}
% The recent advancements in natural language processing literature have seen a significant development of models capable of comprehending and following instructions given in natural language.
Significant progress has been made in the field of natural language processing, leading to the development of models that can comprehend and follow instructions given natural language inputs~\cite{wang2022self, gong2023multimodal, ouyang2022training, brown2020language}. These models harness the power of large language models (LLM) and rely on high-quality instruction data. 
% Similar attempts have been introduced in multi-modal models as well. 
Similarly, efforts have been made to introduce similar capabilities to multi-modal models.
GPT4 \cite{peng2023instruction} has demonstrated impressive performance in multi-modal conversations with humans, yet the techniques contributing to its extraordinary capabilities remain opaque. As a result, several large multi-modal models (LMMs) have recently emerged \cite{zhu2023minigpt, liu2023visual, gong2023multimodal, dai2023instructblip}, such as MiniGPT4 \cite{zhu2023minigpt} and LLaVA \cite{liu2023visual}, both utilize the Vicuna \cite{chiang2023vicuna} as the language generator but with different vision encoders \cite{radford2021learning, li2022blip}. InstructBLIP \cite{dai2023instructblip} is initialized from a pre-trained BLIP-2 \cite{li2023blip} while Multimodal-GPT (MMGPT) \cite{gong2023multimodal} is built on Flamingo \cite{alayrac2022flamingo, anas_awadalla_2023_7733589}.

A recent study~\cite{li2023otter}   revealed that the hallucination issue of LLM, although not desired, is inherited by these LMMs \cite{zhu2023minigpt, liu2023visual, gong2023multimodal, dai2023instructblip}. %Hallucination \cite{bang2023multitask} is one of the major ethical concerns of LLMs, and it can have harmful consequences  
Hallucination, a major ethical concern associated with LLMs~\cite{bang2023multitask}, can lead to harmful consequences,
especially when users without adequate domain knowledge over-rely on these increasingly convincing language models.
In the context of LMM hallucinations, the model can generate descriptions with conflicting information to the given image. For instance, as shown in Fig. \ref{fig:model}~(highlighted in red), existing LMMs \cite{zhu2023minigpt, liu2023visual, dai2023instructblip, gong2023multimodal} tend to describe nonexistent objects such as a "dog" engaging in a nonexisting activity like "playing with a ball".
Additionally, the model may generate long image descriptions without following human instructions (highlighted in blue). 

\textbf{\textit{What are the likely causes of these hallucinations?}} As current LMMs are built on strong LLMs, they may over-rely on language priors and generate words more likely to go together with the instruction text regardless of the image content. What's more, LMMs, such as MiniGPT4 \cite{zhu2023minigpt} and LLaVA \cite{liu2023visual}, employ synthetic instruction data for training,  which are generally long and involve nonexistent objects, activities, or relationships in the image.

\textbf{\textit{Why can't LMMs accurately follow human instructions?}} We conjecture it is due to the lack of diversity in their training data. For example, MiniGPT4 \cite{zhu2023minigpt} is only instruct tuning with four instruction templates designed for image captioning tasks. Though MMGPT \cite{gong2023multimodal} and InstructBLIP \cite{dai2023instructblip} combine several datasets as the instruction tuning data, their instructions and answers are still based on a few templates.

To address these challenges, we present \textit{LRV-Instruction}, a large and diverse visual instruction benchmark. Our benchmark consists of 120k visual instructions generated by GPT4, taking inspiration from the success of recent GPT models in text-annotation tasks~\cite{liu2023gpteval}. Unlike previous studies that focused on limited tasks and pre-defined templates created by human experts~\cite{zhu2023minigpt, dai2023instructblip, gong2023multimodal}, \textit{LRV-Instruction} covers 16 vision-language tasks with open-ended instructions and answers, as shown in Fig.~\ref{fig:tasks} and Tab.~\ref{tab:stat1}. As observed by \cite{li2023evaluating}, current LMMs tend to answer \textit{"Yes"} for any instructions presented to the model, even when the proper answer should be \textit{"No"}. Some models even answer with \textit{"Yes"} for over 90\% of the negative testing examples during their evaluation. Our investigation reveals that most LMMs are finetuned on unbalanced datasets containing only positive instructions (Tab. \ref{tab:comparison}). To enable LMMs to respond to human instructions more faithfully, we design \textit{LRV-Instruction} to include both negative and positive instructions for robust instruction tuning. Our negative instructions are generated at two semantic levels (Fig. \ref{fig:neg_example}): \textit{(i) Nonexistent Element Manipulation} and \textit{(ii) Existent Element Manipulation} in two different formats, \textit{Declarative} and \textit{Interrogative}. To improve the robustness and flexibility of the evaluation on visual instruction tuning, we propose \textit{GPT4-Assisted Visual Instruction Evaluation (GAVIE)} to assess the LMM output in two different aspects: \textit{Relevancy} to evaluate the instruction-following performance and \textit{Accuracy} to measure the visual hallucination in the LMM output. \textit{GAVIE} does not require human-annotated groundtruth answers \cite{rohrbach2018object} and can be easily adapted to different instruction formats instead of specific designs in \cite{li2023evaluating}.

% Based on our experiments,
We empirically evaluate four publicly available LMMs~\cite{zhu2023minigpt,liu2023visual,dai2023instructblip,gong2023multimodal} on our benchmark and found that existing LMMs seriously hallucinate when prompted with our negative instructions. We further verify the effectiveness of our \textit{LRV-Instruction} by finetuning MiniGPT4~\cite{zhu2023minigpt} on this more balanced data. Our instruct-tuned MiniGPT4 suffers much less from hallucination and achieves state-of-the-art performance compared to the original MiniGPT4, LLaVA \cite{liu2023visual}, InstructBLP \cite{dai2023instructblip} and MMGPT \cite{gong2023multimodal} on both our evaluation set and public benchmarks \cite{li2023evaluating, hudson2019gqa}. We also observe that \textit{Existent Element Manipulation} instructions are more challenging than \textit{Nonexistent Element Manipulation} instructions for LMMs. Furthermore, a robust model performance requires a balanced ratio between positive and negative instances. To sum up, our contributions are three-fold:
\begin{itemize}[nosep,leftmargin=*]
    \item We build \textit{LRV-Instruction}, a large and diverse dataset containing 120k visual instructions, with 16 vision and language tasks and negative instructions in different semantic levels and styles.
    \item We propose \textit{GAVIE}, a novel approach to evaluate visual instruction tuning without requiring groundtruth answers and pre-designed instruction formats.
    \item  We conduct comprehensive experiments to investigate the hallucination of current LMMs. The empirical study validates the effectiveness of \textit{LRV-Instruction} for robust visual instruction tuning.
\end{itemize}

\begin{figure*}[t]
    \centering
      \includegraphics[width=1\textwidth]{task2.pdf}
    \caption{Examples of negative instructions in our \textit{LRV-Instruction} dataset and output from current LMMs and our model. \textbf{\textcolor{red}{RED}} means inconsistent elements in the negative instructions. }
    \label{fig:neg_example}
    \vspace{-0.15in}
\end{figure*}




\label{sec:Introduction}

\section{Related Works}
Early explorations \cite{wang2022git, li2022blip,li2019visualbert, sun2019videobert} of vision and language pre-trained models tend to use Bert-based \cite{liu2019roberta, koroteev2021bert} models as the language decoder. Inspired by the recent success of large language models \cite{touvron2023llama, gilardi2023chatgpt, zhao2023survey} and datasets \cite{lin2014microsoft, changpinyo2021conceptual, krishna2017visual, liu2020visual, sharma2018conceptual, srinivasan2021wit, liu2023documentclip}, many studies \cite{alayrac2022flamingo, li2023blip, li2023otter} have been focused on improving vision-language pre-trained models by integrating powerful LLMs with in-context or few-shot learning capability. More recently, some visual instruction-tuned LMMs \cite{zhu2023minigpt, liu2023visual, gong2023multimodal, dai2023instructblip} have emerged, showing excellent generalization performance in unseen VL tasks. Specifically, LLaVA \cite{liu2023visual} projects the output of a visual encoder as input to LLaMA \cite{touvron2023llama} and trains both the alignment network and the LLM on synthetic data. MiniGPT4 \cite{zhu2023minigpt} is built on BLIP-2 but uses Vicuna \cite{chiang2023vicuna} as the language decoder. It only finetunes the cross-modal alignment network on longer image captions from ChatGPT. The research approaches \cite{gong2023multimodal, dai2023instructblip} are instruction-tuned on a collection of VL datasets, but InstructBLIP \cite{dai2023instructblip} uses BLIP2 \cite{li2023blip} as the backbone while \cite{gong2023multimodal} is initialized from Flamingo \cite{alayrac2022flamingo}. mPLUG-owl \cite{ye2023mplug} finetunes LLaMA \cite{touvron2023llama} model using both text instruction data and vision-language instruction data from LLaVA \cite{liu2023visual}. In comparison, we propose a large and diverse visual instruction dataset with 16 vision and language tasks and negative instructions in different semantic levels and styles. This can help improve the robustness of current LMMs.

Although LMMs are powerful in solving VL tasks, they also suffer from the hallucination  inherited from LLM. Popular image captioning metrics like CIDEr \cite{vedantam2015cider} SPICE \cite{anderson2016spice} do not appropriately penalize hallucination. CHAIR, \cite{rohrbach2018object}, is unstable and needs complex human-crafted parsing rules for exact matching. Alternatively, \cite{li2023evaluating} converts the hallucination into a binary classification problem. However, it requires the input questions to follow specific templates, such as \textit{"Is there a/an <object> in the image?"}. In comparison, our proposed GAVIE  can evaluate model hallucination in an  open-ended manner without needing human-annotated groundtruth answers.


\begin{figure*}[t]
    \centering
      \includegraphics[width=1\textwidth]{tasks.pdf}
    \caption{Examples of positive instructions in our \textit{LRV-Instruction} dataset. Each instance shows the image with the corresponding task and instruction. More examples are in the Appendix.}
    \label{fig:tasks}
\vspace{-0.17in}
\end{figure*}
\label{sec:related_work}


% \section{Dataset}
\section{LRV-Instruction}
\subsection{Automatic Visual Instruction Generation}
Annotating large-scale visual instruction data can be challenging and time-consuming \cite{wang2022self}. 
 It involves expertly written detailed instructions and specific labels for different tasks. Inspired by the success of GPT4 in text-annotation tasks \cite{gilardi2023chatgpt}, we leverage GPT4, instead of human  workers, to build \textit{LRV-Instruction}. \textit{LRV-Instruction} is designed to cover a variety of VL tasks (Fig. \ref{fig:tasks}), with open-ended positive and negative instructions (Fig. \ref{fig:neg_example}, \ref{fig:tasks}) in different linguistic styles (Tab.\ref{tab:stat2}).

In what follows, we detail the construction process.

\textbf{Prompt Design.} Inspired by \cite{wang2022self}, we use the in-context few-shot learning ability of GPT4 to automatically generate instruction data for various VL tasks. We filter the output tasks manually and select 16 tasks (Tab. \ref{tab:stat1}) with text answers.  In contrast with  \cite{liu2023visual} using a few scene captions to represent an image as input to the text-only GPT4, we take advantage of the Visual Genome dataset \cite{krishna2017visual}, which has detailed visual information like image size, bounding boxes, and dense captions. Specifically, each image typically has 21 object regions and their corresponding captions. We leverage GPT4 % as the teacher 
to create the instruction-following data with the image size, bounding boxes, and dense captions as the "visual" input as if it can "see" the image. An example is shown in Fig. \ref{fig:prompt1}. There can be more than one caption describing the same object with different attributes (Fig. \ref{fig:prompt1}), such as "\textit{woman wearing a long dress}" and "\textit{woman wearing a yellow dress}". Although we present the bounding box coordinates of each caption to GPT4, it can be easily confused, treating them as two instances, one in a long dress and the other in a yellow dress.
% according to the bounding box coordinates
To mitigate this issue, we add "\textit{highly overlapping bounding boxes may refer to the same object}" into the prompt to help GPT4 understand the "visual" input better. For each image, we randomly select 10 tasks. 
%In order to enhance the diversity of our dataset,
To enrich the instructions, we ask GPT4 to generate instances in both declarative and interrogative formats. The limitation of \cite{liu2023visual, zhu2023minigpt} is that synthetic visual instructions are generally longer and may involve unexpected descriptive information that is inconsistent with the image. Therefore, we explicitly instruct GPT4 with \textit{"The answers should be less than 30 words"} as a requirement to reduce the chance of generating extra unrelated information in the training data, which can also improve the reading experience for users.

\textbf{Negative Instruction Generation. }As shown in \cite{li2023evaluating}, current LMMs tend to answer \textit{“Yes”} by following any instruction presented to the model rather than predicting a faithful answer. 
% honestly saying \textit{"no"}. 
To teach LMMs \cite{zhu2023minigpt, liu2023visual, gong2023multimodal, dai2023instructblip} to answer questions in instructions faithfully, we introduce two categories of negative instructions: \textit{(1) Neg1: "Nonexistent Element Manipulation"} by introducing nonexistent objects, activities, attributes and interactions to the "visual" input as described above.
% the image. 
\textit{(2) Neg2: "Existent Element Manipulation"} by manipulating existent objects with inconsistent attributes (Fig. \ref{fig:neg_example}). As for the detailed prompt of \textit{Neg1}, we leverage the same format of the "visual" input as shown in Fig. \ref{fig:prompt1}. Additionally, we provide the following instructions to GPT4: 

\textit{"Come up with 6 misleading instructions with \underline{nonexistent elements (nonexistent objects, nonexistent} \underline{activities, nonexistent attributes, nonexistent interactions)} in the images with different language styles. The instructions should contain interrogative and declarative sentences. Please also explain the reason."}

\begin{table}[t]
\setlength\tabcolsep{4.3pt}
\centering
\small
\begin{tabular}{lccccc}
\toprule[1.5pt]
 & Ours & MiniGPT4 & LLaVA & InstructBLIP & MMGPT \\
\midrule
Hard Negative Instructions?      &  \textcolor{green}{\ding{52}} &  \textcolor{red}{\ding{56}} &  \textcolor{red}{\ding{56}} &  \textcolor{red}{\ding{56}} &  \textcolor{red}{\ding{56}}\\
Self Generated Instruction?      &  \textcolor{green}{\ding{52}} &  \textcolor{red}{\ding{56}} &  \textcolor{green}{\ding{52}} &  \textcolor{red}{\ding{56}} &  \textcolor{red}{\ding{56}}\\
Address Hallucination?      &  \textcolor{green}{\ding{52}} &  \textcolor{red}{\ding{56}} &  \textcolor{red}{\ding{56}} &  \textcolor{red}{\ding{56}} &  \textcolor{red}{\ding{56}}\\
NOT Template Instruction?    &  \textcolor{green}{\ding{52}} &  \textcolor{red}{\ding{56}} &  \textcolor{green}{\ding{52}} &  \textcolor{red}{\ding{56}} &  \textcolor{red}{\ding{56}}\\
\# of Instances      &  120k &  3k &  150k &  >1M &  >500k\\
\# of VL Tasks      &  16 &  1 &  3 &  11 &  5\\
\bottomrule[1.5pt]
\end{tabular}
\vspace{0.05in}
\caption{
A comparison of \textit{LRV-Instruction} with datasets used by current LMMs. InstructBLIP and MMGPT just combine existing VQA datasets while we create a new dataset by GPT4.}
\label{tab:comparison}
% \vspace{-0.28in}
\vspace{-6.5mm}
\end{table}

\begin{figure*}[t]
    \centering
      \includegraphics[width=1
\textwidth]{prompt11.pdf}
    \caption{One example to illustrate the prompt we use to generate the visual instruction data by GPT4. We use the bounding box coordinates and dense captions to represent image content. More examples can be found in the Appendix.}
    \label{fig:prompt1}
\vspace{-4mm}
\end{figure*}


We replace the underlined text above with \textit{"existing objects but wrong attributes"} for the prompt of \textit{Neg2}. We present the complete prompts and example instructions in the Appendix. % In the appendix, we present the complete prompts and example results in Fig. \ref{fig:prompt_demo1}, \ref{fig:prompt_demo2}, \ref{fig:demo_neg1} and \ref{fig:demo_neg2}. % show more instruction demos in the appendix.

\textbf{Quality Control. }
% To ensure quality so that 
% To ensure the quality of the training data, we exclude instructions employing dense captions from the visual input as default answer choices for the multi-choice VQA task. \kevin{this statement is not clear}
For multi-choice VQA tasks, we observe that the generated instruction may occasionally contain no choices, and we remove such incomplete data.
We also remove instances with answers longer than 30 words. To examine the quality of our dataset, we randomly sample 500 instances and ask ten expert annotators to determine whether the output answers from GPT4 are correct or not, with regard to the instruction and the image content. We found 91\% of the instructions are appropriate for the image inputs. Furthermore, 85\% of outputs are acceptable responses to the instructions. Even though some responses may contain errors, most generations conform to the correct structure, serving as applicable visual instruction tuning guidelines.


\subsection{Data Statistics}
Tab. \ref{tab:stat2} shows the statistics of the generated data. We created a total of over 120k visual instructions after filtering. Tab. \ref{tab:comparison} shows a comparison of \textit{LRV-Instruction} and other datasets used by current LMMs. \textit{LRV-Instruction} covers much more VL tasks than existing visual instruction tuning datasets. Instead of only using positive instructions, \textit{LRV-Instruction} also includes negative instructions at different semantic levels for robust instruction tuning. In addition, employing the GPT4-assisted generation, \textit{LRV-Instruction} has more open-ended instructions instead of following a few templates. Compared to relevant datasets, these comparisons demonstrate the advantages of \textit{LRV-Instruction}.



\label{sec:dataset}

\section{Visual Instruction Tuning}
\vspace{-0.1in}
% In this paper, our model architecture is initialized from 
We constructed a model for visual instruction tuning based on MiniGPT4 \cite{zhu2023minigpt}, which consists of the Vision transformer \cite{liu2023covid} backbone as the image encoder, Vicuna \cite{chiang2023vicuna} as the text decoder and a pre-trained Q-Former to connect them. Vicuna is built upon LLaMA \cite{touvron2023llama} with stronger following ability. Following \cite{zhu2023minigpt}, the Q-Former is designed to extract visual features from the frozen image encoder. Before feeding into the frozen Vicuna as the visual prompt, we use a learnable linear projection layer to narrow the gap between extracted visual features with Vicuna embeddings. The detailed prompt is as follows: 

\textit{"Give the following image: <Img>ImageContent</Img>. You will be able to see the image once I provide it to you. Please answer my questions. \#\#\#Human: <Img><ImageHere></Img> Instruction \#\#\#Assistant: Answer"},

where \textit{"<Img>ImageContent</Img>"} is a template to let Vicuna know image features will be placed within the two special tokens \textit{"<Img>", "</Img>"}. \textit{"<ImageHere>"} is the actual slot for image features. \textit{"Instruction"} and \textit{"Answer"} represent the human instructions and model output text, respectively.
\label{sec:method}

\section{Evaluation Setting}
\vspace{-0.1in}
\begin{table}[t]
\begin{subtable}[t]{0.3\textwidth}
\small
\begin{tabular}{c}
\toprule[1.5pt]      
Tasks of \textit{LRV-Instruction} \\ 
\midrule
\midrule  
\pbox{7cm}{\textit{Image Captioning, Object Detection, Image Sentiment Analysis, Image Quality Assessment, Object Interaction Analysis, Image Anomaly Detection, Referential Expression Grounding, OCR, VCR, Object Attribute Detection, Muli-choice VQA, Semantic Segmentation, Dense Captioning, Visual Entailment, Styled Image Caption, Activity Recognition, Facial Expression Detection
}}\\
\bottomrule[1.5pt]
\end{tabular}
 \caption{VL tasks in our dataset.}
\label{tab:stat1}
\end{subtable}
\hspace{\fill}
\begin{subtable}[t]{0.5\textwidth}
\setlength\tabcolsep{4.7pt}
\centering
\small
\begin{tabular}{lc}
\toprule[1.5pt]
Statistic &\\
\midrule
\# of Instructions      & 120k\\
\# of Images         & 6.5k\\
\# of Interrogative Instructions        & 70k \\
\# of Declarative Instructions        & 50k \\
Ave \# of Instruction per Image       & 17.7 \\
Ave Answer Length      & 15.2 \\
Ave Instruction Length     & 12.1 \\
\midrule
\# of Negative Instructions     & 35k\\
Ave Instruction Length (Neg1) & 18.1\\
Ave Instruction Length (Neg2) & 11.2\\
\% of Interrogative (Neg1) & 47\% \\
\% of Interrogative (Neg2) & 81\% \\
\bottomrule[1.5pt]
\end{tabular}
\vspace{0.05in}
\caption{
Statistics of \textit{LRV-Instruction}.
}
\label{tab:stat2}
\end{subtable}
\caption{ (a) The proposed \textit{LRV-Instruction} covers a diverse set of vision-language tasks. (b) Detailed statistics of our constructed dataset. 
Neg1 means \textit{Nonexistent Element Manipulation} and Neg2 means \textit{Existent Element Manipulation} in the negative instructions.
}
\label{tab:stat}
\vspace{-0.2in}
\end{table}

CHAIR \cite{rohrbach2018object} was introduced to evaluate object hallucination in image captioning tasks. However, it 
usually demands complex human-crafted rules. Alternatively, \cite{li2023evaluating} formulates the evaluation of object hallucination as a binary classification task that prompts LMM to output "Yes" or "No". However, it is hard to evaluate the LMM output in an open-ended manner. In addition, both methods highly depend on human-annotated groundtruth answers.


To this end, we introduce \textit{GPT4-Assisted Visual Instruction Evaluation (GAVIE)} as a more flexible and robust approach. The general prompt we use is shown in %Fig. \ref{fig:prompt_eval} in 
the Appendix. GPT4 takes the dense captions with bounding box coordinates as the image content and compares human instructions and model response. Then we ask GPT4 to work as a smart teacher and score (0-10) students' answers based on two criteria. \textit{(1) Accuracy: whether the response is accurate concerning the image content}. \textit{(2) Relevancy: whether the response directly follows the instruction}. An evaluation result is shown in Fig. \ref{fig:model}, which successfully points out that \textit{"dog, ball"} is inconsistent with the image, and the response from the first LMM did not address the instruction. Unlike previous evaluation methods \cite{liu2023visual, li2023evaluating, rohrbach2018object}, \textit{GAVIE} does not require human-annotated groundtruth answers and can freely adapt to diverse instruction formats.

\begin{table}[t]
\setlength\tabcolsep{3pt}
\centering
\small
\begin{tabular}{lccccc}
\toprule[1.5pt]
& Ours & MiniGPT4 & LLaVA & InstructBLIP & MMGPT\\
\midrule
\textsc{Accuracy(GPT4)} & \textbf{6.58} & 4.14 & 4.36 & 5.93 & 0.91\\
\textsc{Relevancy(GPT4)} & \textbf{8.46} & 5.81 & 6.11 & 7.34 & 1.79\\
\bottomrule[1.5pt]
\end{tabular}
\vspace{0.05in}
\caption{
Comparison Experiment on our evaluation set evaluated by \textit{GAVIE}. The metric scores of \textsc{Accuracy(GPT4)} and \textsc{Relevancy(GPT4)} are in a scale of 0 to 10.
}
\label{tab:comparison_evaluation}
\vspace{-0.1in}
\end{table}


\begin{table}[t]
\setlength\tabcolsep{3pt}
\centering
\small
\begin{tabular}{lcccc}
\toprule[1.5pt]
Model & Accuracy & F1 Score\\
\midrule
LLaVA & 0.45 & 0.59 \\
% \midrule
MiniGPT4 & 0.44 & 0.52\\
% \midrule
InstructBLIP & 0.74 & 0.81 \\
% \midrule
Ours & \textbf{0.79} & \textbf{0.84} \\
\bottomrule[1.5pt]
\end{tabular}
\vspace{0.05in}
\caption{Comparison experiment on hallucination evaluation dataset introduced by \cite{li2023evaluating}.} 
\label{tab:coco}
\vspace{-0.2in}
\end{table}






\label{sec:Evaluation}

\section{Experiment}
\vspace{-0.1in}
\subsection{Implementation Setup}
\textbf{Baselines. }We evaluate the zero-shot performance of 4 recently released LMMs: (1) MiniGPT4 (7B) \cite{zhu2023minigpt}; (2) LLaVA (7B) \cite{liu2023visual}; (3) InstructBLIP (7B) \cite{dai2023instructblip}; (4) Multimodal-GPT (7B) (MMGPT) \cite{gong2023multimodal}. All the above models have been tuned on their own collected visual instruction data. 


\textbf{Training Details. }We initialize the model from the first pretraining stage checkpoint of MiniGPT4. Then we instruct-tune the model on our dataset with the linear projection layer as the only learnable module. As for the hyper-parameters, we  employ a batch size of 1, a linear warmup of the learning rate 1e-6 during the initial 200 steps, and a weight decay of weight decay 0.05. We trained our models on NVIDIA Quadro RTX 8000 %8 Nvidia A100 (48G) 
GPUs with the maximum epoch as 20. 
% In addition, 
We randomly select 1000 instances from our dataset for evaluation, and the others are used for finetuning.


% \kevin{The below question list can be removed. Repeated in Sec6.2 already}

% \textbf{Research Questions. }We conduct empirical experiments to investigate the following questions: 
% \begin{itemize}[nosep,leftmargin=*]
%     \item \textit{Question1: "Do LMMs perform better on Positive or Negative instructions?"}. 
%     \item \textit{Question2: "Do LMMs perform better on different formats and lengths of instructions?"}. 
%     \item \textit{Question3: "How do LMMs perform on different semantic levels of hallucination?"}
%     \item \textit{Question4: "How do LMMs perform on different datasets?"}
%     \item \textit{Question5: "How do LMMs perform on different ratios of composition in training data?"} 
% \end{itemize}


\subsection{Main Results}
% \textbf{Results on \textit{LRV-Instruction}. }
\textbf{\textit{How do LMMs perform on \textit{LRV-Instruction}? }}We show the evaluation results on our dataset in Tab.~\ref{tab:comparison_evaluation}. Among the baselines, InstructBLIP achieves better results than other LMM baselines because its visual instructions are collected from a wide variety of publicly available datasets. LLaVA \cite{liu2023visual} utilizes the GPT-assisted approach to generate visual instructions, but its performance is much worse. This is probably because its synthetic answers from GPT4 are generally longer and may involve irrelevant information. As a comparison, our model outperforms the existing LMM baselines by a large margin, which benefits from the rich composition of our dataset and better prompt design. %\kevin{I would suggest to have a subsection "Main Result" at the beginning of Sec 6.2, and put this paragraph in the subsection. After that, add the title "How do LMMs perform on different datasets?" to the next paragraph. }
%\kevin{Another way: Change "How do LMMs perform on different datasets?" to "Effectiveness of LRV-Instruct..."} \kevin{Another way: at the beginning, mention we test on three datasets, including AAA,BBB,CCC. Then discuss each result.}


% \begin{table}[t]
% \begin{subtable}[t]{0.5\textwidth}
% \centering
% \small
% \begin{tabular}{lcc}
% \toprule[1.5pt]
% Model & Training Size & Accuracy\\
% \midrule
% InstructBLIP & >1M & \textbf{0.62}\\
% LLaVA & 150k & 0.48\\
% MiniGPT4 & 3k & 0.44\\
% Ours & 120k& \textbf{0.61}\\
% \bottomrule[1.5pt]
% \multicolumn{1}{c}{} \\ % space holder
% \end{tabular}
%  \caption{Evaluation results on GQA dataset \cite{hudson2019gqa}.}
% \label{tab:perform11}
% \end{subtable}
% \hspace{\fill}
% \begin{subtable}[t]{0.5\textwidth}
% \centering
% \small
% \begin{tabular}{lcc}
% \toprule[1.5pt]
% Ratio & \textit{$Acc_p$} & \textit{$Acc_n$}\\
% \midrule
% \textit{All Pos} & 0.97 & 0.05\\
% \textit{Pos:Neg=2:1} & 0.95 & 0.50\\
% \textit{Pos:Neg=1:1} & 0.92 & 0.85\\
% \textit{Pos:Neg=1:2} & 0.87 & 0.86\\
% \textit{All Neg} & 0.10 & 0.98\\
% \bottomrule[1.5pt]
% \end{tabular}
% \caption{
% Results of different composition ratios in training.
% }
% \label{tab:perform112}
% \end{subtable}
% \vspace{-0.05in}
% \caption{
% Evaluation of current LMMs on the VQA dataset and ablation study of our model.
% }
% \label{tab:perform12}
% \vspace{-0.2in}
% \end{table}

%%%%%%%

\begin{table}[t]
\centering
\def\arraystretch{1.2}
\begin{minipage}{0.48\linewidth}\centering
\centering
\small
\begin{tabular}{lcc}
\toprule[1.5pt]
Model & Training Size & Accuracy\\
\midrule
InstructBLIP & >1M & \textbf{0.62}\\
LLaVA & 150k & 0.48\\
MiniGPT4 & 3k & 0.44\\
Ours & 120k& \textbf{0.61}\\
\bottomrule[1.5pt]
\multicolumn{1}{c}{} \\ % space holder
\end{tabular}
\vspace{-0.05in}
\caption{Results on GQA~\cite{hudson2019gqa}. We evaluate LMM's performance in general VQA scenario.}
\label{tab:perform11}
\end{minipage}
%
\begin{minipage}{0.48\linewidth}\centering
\begin{tabular}{lcc}
\toprule[1.5pt]
Ratio & \textit{$Acc_p$} & \textit{$Acc_n$}\\
\midrule
\textit{All Pos} & 0.97 & 0.05\\
\textit{Pos:Neg=2:1} & 0.95 & 0.50\\
\textit{Pos:Neg=1:1} & 0.92 & 0.85\\
\textit{Pos:Neg=1:2} & 0.87 & 0.86\\
\textit{All Neg} & 0.10 & 0.98\\
\bottomrule[1.5pt]
\end{tabular}
\vspace{0.1in}
\caption{Results of different composition ratios in visual instruction tuning. }
\label{tab:perform112}
\end{minipage}
\end{table}



% \begin{table*}[t!]\centering
% %%%%%%%%%% subfloat a
% \subfloat[Evaluation results on GQA dataset \cite{hudson2019gqa}\label{tab:perform11}]
% \begin{tabular}{lcc}
% 	\toprule[1.5pt]
% 	Model & Training Size & Accuracy\\
% 	\midrule
% 	InstructBLIP & >1M & \textbf{0.62}\\
% 	LLaVA & 150k & 0.48\\
% 	MiniGPT4 & 3k & 0.44\\
% 	Ours & 120k& \textbf{0.61}\\
% 	\bottomrule[1.5pt]
% \multicolumn{1}{c}{} \\ % space holder
% \end{tabular}
% \hfill
% %%%%%%%%%% subfloat b
% \subfloat[Results of different composition ratios in training \label{tab:perform112}]
% \begin{tabular}{lcc}
% 	\toprule[1.5pt]
% 	Ratio & \textit{$Acc_p$} & \textit{$Acc_n$}\\
% 	\midrule
% 	\textit{All Pos} & 0.97 & 0.05\\
% 	\textit{Pos:Neg=2:1} & 0.95 & 0.50\\
% 	\textit{Pos:Neg=1:1} & 0.92 & 0.85\\
% 	\textit{Pos:Neg=1:2} & 0.87 & 0.86\\
% 	\textit{All Neg} & 0.10 & 0.98\\
% 	\bottomrule[1.5pt]
% \end{tabular}
% \caption{
% Evaluation of current LMMs on the VQA dataset and ablation study of our model
% }
% \vspace{-4mm}
% \end{table*}



\begin{table}[t]
\setlength\tabcolsep{3pt}
\centering
\small
\begin{tabular}{lcccccc}
\toprule[1.5pt]
Categories& Metric&Ours & MiniGPT4 & LLaVA & InstructBLIP & MMGPT\\
\midrule
Neg1 & \textsc{Accuracy(GPT4)} & \textbf{8.90} & 3.72 & 2.09 & 5.50 & 1.13\\
Neg1 & \textsc{Relevancy(GPT4)} & \textbf{8.96} & 5.94 & 4.83 & 7.22 & 2.24\\
\midrule
Neg2 & \textsc{Accuracy(GPT4)} & \textbf{6.50} & 2.57 & 1.42 & 2.18 & 0.96\\
Neg2 & \textsc{Relevancy(GPT4)} & \textbf{8.46} & 2.53 & 1.82 & 2.73 & 1.19\\
\bottomrule[1.5pt]
\end{tabular}
\vspace{0.05in}
\caption{
Evaluation results on \textit{Neg1} instructions and \textit{Neg2} instructions by \textit{GAVIE}. The metric scores of \textsc{Accuracy(GPT4)} and \textsc{Relevancy(GPT4)} are in a scale of 0 to 10.}
\label{tab:neg1neg2}
\vspace{-0.1in}
\end{table}


\textbf{\textit{How do LMMs perform on different datasets? }}We compare our model against the baseline models on the hallucination evaluation dataset recently introduced by POPE~\cite{li2023evaluating}, which randomly selects 500 images from COCO \cite{lin2014microsoft}. Since the authors of POPE~\cite{li2023evaluating} have not released their dataset at the time of our paper submission, we employ our GPT4-assisted approach to generate open-ended negative instructions. In addition, we instruct the model to produce “Yes” or “No” and use classification accuracy and F1 for evaluation according to \cite{li2023evaluating}. As shown in Tab. \ref{tab:coco}, we observe that MiniGPT4, LLaVA and InstructBLIP may not work well with open-ended negative instructions. In contrast, the highest scores of our model demonstrates that our dataset is more suitable for robust visual instruction tuning. Although InstructBLIP tends to hallucinate less, its training set is much larger than ours (\textit{120k vs 1M}). Moreover, our model is built on MiniGPT4, which only instruct-tunes the linear projector layer. To further explore the LMMs' performance in the common scenario of visual question-answering (VQA), we follow \cite{li2023evaluating} to evaluate LMMs on GQA dataset~\cite{hudson2019gqa} in the zero-shot setting. It is worth noting that GQA consists of open-ended VQA questions. As shown in Tab. \ref{tab:perform11}, the results suggest that our method achieves on-par performance with InstructBLIP in generic VQA setting, and our method requires much less training data. %By finetuning MiniGPT4 on our dataset, it achieves on-par performance with InstructBLIP but employs much less training data.


%we follow their proposed method to construct the \textsc{Random Negative} instructions. To construct suchinstructions, we randomly sample the objects from 80 object classes in COCO that do not exist in the image. In addition, we instruct the model to produce “Yes” or “No” and use classification accuracy and F1 for evaluation. We adopt the same instruction templates from \cite{li2023evaluating}: \textit{(1) Is there an <object> in the image? (2) Does the image contain a <object>?} 

%As shown in Tab. \ref{tab:coco}, our model outperforms MiniGPT4 and LLaVA in the \textsc{Random Negative} instructions which demonstrates the effectiveness of our dataset in mitigating the hallucination issue. Although InstructBLIP tends to hallucinate less, its training set is much larger than ours (\textit{120k vs 1M}). Moreover, our model is built on MiniGPT4, which only instruct-tunes the linear projector layer. 

%In order to evaluate LMMs on non-template instructions in this dataset, we employ our GPT4-assisted approach to generate \textsc{Hard Negative} instructions. 

%From Tab.~\ref{tab:coco}, the overall performance of baselines is worse compared to the \textsc{Random Negative} setting, showing that baselines may not work well with open-ended negative instructions. In contrast, our model can achieve similar performance for both negative instruction categories, which demonstrates that our dataset is more suitable for robust visual instruction tuning. 





\subsection{Detailed Analysis}
\textbf{\textit{Do LMMs perform better on Positive or Negative Instructions? }}Our evaluation set consists of positive and negative instances. We divide it into two sets and analyze the model performance on each of them. As shown in Fig.~\ref{fig:posneg}, baseline models, including MiniGPT4, LLaVa, and InstructBLIP,  perform better on positive instances than negative ones, as the training data adopted by these models do not contain negative instructions. MMGPT has low performance on both sets due to many repetitive phrases in the response. In addition, we found that the degradation of LLaVA is the most severe. We hypothesize that the synthetic answers for instruction tuning in LLaVA are generally longer and involve more unrelated information. In contrast, our model performs the best in both sets. InstructBLIP performs with higher scores than other LMMs because of the effectiveness of its instruction-aware visual encoder to extract image information.



\textbf{\textit{How do LMMs perform at the different semantic levels of hallucination? }} As shown in Tab \ref{tab:neg1neg2}, all baselines perform better on \textit{Neg1} than \textit{Neg2}. From the visual perspective, existent object manipulations with wrong attributes in \textit{Neg2} are more challenging than adding nonexistent objects from images to instructions in \textit{Neg1}. For example, in Tab. \ref{fig:neg_example}, it %is not hard 
may be straightforward to find that the "hot air balloon" does not appear in the image. However, "woman" does exist in the second example of Fig.~\ref{fig:posneg} while she is not in the blue pants and pink shirts, which requires a fine-grained understanding of the visual content. Therefore, a more powerful vision encoder is needed for future LMMs. In contrast, our model improves a lot at both semantic levels, which benefits from our diverse instruction tuning data. In the Appendix, we present more visualizations of model predictions for \textit{Neg1} and  \textit{Neg2} examples. % Fig.~\ref{fig:demo_neg1} and \ref{fig:demo_neg2} in the appendix show more visualizations of model predictions for \textit{Neg1} and  \textit{Neg2} examples. %\linjie{Will Appendix be in the same file as the main text? If not, need to mention Figure \ref{fig:demo_neg1} and \ref{fig:demo_neg2} are in the appendix.}
% of different hallucinations.

% \kevin{we have discussed length and format in the previous paragraph. Here we should directly discuss observations at different semantic levels}




\textbf{\textit{How do LMMs perform at the different composition ratios in training data?}} Our model benefits from the datasets with both positive and negative instructions. In Tab. \ref{tab:perform112}, we investigate how it addresses hallucination issues with different ratios of positive and negative samples in the training set. Inspired by \cite{li2023evaluating}, we instruct the model to produce “Yes” or “No” and use classification accuracy for evaluation. \textit{$Acc_p$} is the accuracy on the positive instruction set, while \textit{$Acc_n$} is the accuracy on the negative instruction set. From Tab. \ref{tab:perform112}, we found that \textit{$Acc_n$} increases with more negative samples, which verifies our hypothesis that the hallucination problem of current LMMs \cite{zhu2023minigpt, liu2023visual, gong2023multimodal, dai2023instructblip} are due to lack of 
negative instructions. Besides, with a balanced ratio (\textit{pos:neg=1:1}), the model performs the best in both positive and negative sets. We also found that the setting \textit{pos:neg=1:2} can achieve a similar level of results as \textit{pos:neg=1:1}. This is probably because the backbone of our model (MiniGPT4) was trained on a large amount of image-caption pair datasets, which are positive samples.


\textbf{\textit{Do LMMs perform better on different formats and lengths of instructions? }}From Tab \ref{tab:Interrogative}, LMMs perform with higher scores on interrogative instructions than declarative, but the difference is relatively small. Even though recent visual instruction tuning datasets lack diverse declarative instructions, the LMMs built on LLM are powerful enough to understand and follow the declarative instructions. From Fig. \ref{fig:length}, current LMMs achieve better results in short instructions than long ones since longer instructions contain more information, making it more challenging to understand. 


\begin{figure}[t]
     \centering
     \begin{subfigure}[b]{0.49\textwidth}
         \centering
         \includegraphics[width=\textwidth]{sub1.pdf}
         \vspace{-0.2in}
         \caption{Accuracy Performance.}
         \label{fig:y equals x}
     \end{subfigure}
     \hfill
     \begin{subfigure}[b]{0.49\textwidth}
         \centering
         \includegraphics[width=\textwidth]{sub2.pdf}
         \vspace{-0.2in}
         \caption{Relevancy Performance.}
         \label{fig:three sin x}
     \end{subfigure}
        \caption{Evaluation results on positive and negative instructions by \textit{GAVIE}. }
        \label{fig:posneg}
\vspace{-0.1in}
\end{figure}

\begin{figure}[t]
     \centering
     \begin{subfigure}[b]{0.49\textwidth}
         \centering
         \includegraphics[width=\textwidth]{length1.pdf}
         \vspace{-0.2in}
         \caption{Accuracy Performance.}
         \label{length1}
     \end{subfigure}
     \hfill
     \begin{subfigure}[b]{0.49\textwidth}
         \centering
         \includegraphics[width=\textwidth]{length2.pdf}
         \vspace{-0.2in}
         \caption{Relevancy Performance.}
         \label{length2}
     \end{subfigure}
        \caption{Evaluation results on different instruction lengths by \textit{GAVIE}. }
        \label{fig:length}
\vspace{-0.2in}
\end{figure}


\begin{table}[h]
\centering
\small
\begin{tabular}{lcccccc}
\toprule[1.5pt]
Categories & Metric &Ours & MiniGPT4 & LLaVA & InstructBLIP & MMGPT\\
\midrule
Interrogative&\textsc{Accuracy(GPT4)} & \textbf{6.61} &  4.14 & 4.60 & 5.95 & 1.01\\
Interrogative&\textsc{Relevancy(GPT4)} & \textbf{8.46} &  6.20 & 5.88 & 7.67 & 2.00\\
\midrule
Declarative&\textsc{Accuracy(GPT4)} & \textbf{6.50} & 3.98 & 3.82 & 5.47 & 0.90\\
Declarative&\textsc{Relevancy(GPT4)} & \textbf{8.21} & 5.39 & 5.84 & 6.64 & 1.62\\
\bottomrule[1.5pt]
\end{tabular}
\vspace{0.05in}
\caption{
Evaluation results on Interrogative Instructions and Declarative Instructions by \textit{GAVIE}. The metric scores of \textsc{Accuracy(GPT4)} and \textsc{Relevancy(GPT4)} are in a scale of 0 to 10.
}
\label{tab:Interrogative}
\vspace{-0.2in}
\end{table}
\label{sec:Experiment}

\section{Conclusion}
\vspace{-0.1in}
In this work, we constructed \textit{LRV-Instruction}, a large and diverse dataset containing 120k visual instructions, covering 16 vision and language tasks with both positive and negative instructions in different semantic levels and styles. With \textit{LRV-Instruction}, we  comprehensively investigated the hallucination of existing LMMs and empirically validated its effectiveness in a more robust visual instruction tuning. In addition, we propose \textit{GAVIE}, a novel approach to evaluate visual instruction tuning without requiring human-labeled groundtruth answers and can be easily adapted to different instruction formats. We hope our work can help address the unexpected hallucination issues of LMMs. Future directions include replacing the vision encoders in current LMMs with more powerful visual models to match the capabilities of multimodal GPT4 and investigation of other biases of LMMs to develop more robust models.




\label{sec:Conclusion}


\bibliographystyle{plain}
\bibliography{anthology}


%\bibliographystyle{plain}
%\bibliography{anthology,custom}
%\bibliographystyle{acl_natbib}


\clearpage
\appendix
\section{Appendix}
\subsection{Prompt Design}
\subsubsection{Positive Instance Generation}
We show two full examples of our input prompts in (i) Fig. \ref{fig:prompt_pos11}, \ref{fig:prompt_pos12}, \ref{fig:prompt_pos13} and (ii) Fig. \ref{fig:prompt_pos21}, \ref{fig:prompt_pos22}, \ref{fig:prompt_pos23}. In Fig. \ref{fig:prompt_pos11} and Fig. \ref{fig:prompt_pos21}, we first present the images for the two examples but they are not included in the text prompt for GPT4. As for the text input, we leverage the groundtruth bounding boxes and dense captions to represent the visual content as if GPT4 can see the image. After that, we randomly select 10 tasks from the 16 seeds and ask GPT4 to generate 20 instances for these tasks. Additionally, there can be more than one caption describing the same object with different attributes such as "\textit{woman wearing a long dress}" and "\textit{woman wearing a yellow dress}" in Fig. \ref{fig:prompt_pos11}. Although we present the bounding box coordinates of each caption to GPT4, it can be easily confused, treating them as two instances, one in a long dress and the other in a yellow dress. To mitigate this issue, we add "\textit{highly overlapping bounding boxes may refer to the same object}" into the prompt to help GPT4 understand the "visual" input better. To enrich the instructions, we ask GPT4 to generate instances in both declarative and interrogative formats. We also explicitly instruct GPT4 with \textit{"The answers should be less than 30 words"} as a requirement to reduce the chance of generating extra unrelated information in the training data. In order to make the output of GPT4 in a good format, we also ask GPT4 to generate an instruction, an answer, and a task name in order at the end of the prompt (Fig. \ref{fig:prompt_pos11} and Fig. \ref{fig:prompt_pos21}). The full output of instructions and answers are shown in Fig. \ref{fig:prompt_pos12}, \ref{fig:prompt_pos13} and Fig. \ref{fig:prompt_pos22}, \ref{fig:prompt_pos23}. We also present more positive instances with the output from different LMMs in Fig. \ref{fig:demo_pos1}, \ref{fig:demo_pos2}, \ref{fig:demo_pos3}.




\subsubsection{Negative Instance Generation}
We show two full examples of our input prompts in (i) Fig. \ref{fig:prompt_neg11}, \ref{fig:prompt_neg12} and (ii) Fig. \ref{fig:prompt_neg21}, \ref{fig:prompt_neg22}. In Fig. \ref{fig:prompt_neg11} and Fig. \ref{fig:prompt_neg21}, we present the images to help readers understand dense captions better but they are not included in the text prompt for GPT4. We leverage the bounding boxes and dense captions as the "visual" input. As for  \textit{Nonexistent Element Manipulation} in \ref{fig:prompt_neg11}, we ask GPT4 to generate 6 instructions with nonexistent elements (nonexistent objects, nonexistent activities, nonexistent attributes, nonexistent interactions). As for  \textit{Existent Element Manipulation} in \ref{fig:prompt_neg21}, we ask GPT4 to generate 6 instructions of existing objects with wrong attributes. At the end of the text prompt, we ask GPT4 to generate an instruction and a reason to explain why the instruction is inconsistent with the image in order. The reason is regarded as the answer for the instruction in our training data. Fig. \ref{fig:prompt_neg12} and Fig. \ref{fig:prompt_neg22} show the full output from GPT4. We also present more negative instances with the output from different LMMs in Fig. \ref{fig:demo_neg1}, \ref{fig:demo_neg2}.
.
\subsubsection{GAVIE Evaluation}
We show two full examples of the text prompt for GAVIE in (i) Fig. \ref{fig:prompt_neg31}, \ref{fig:prompt_neg32}, \ref{fig:prompt_neg33} and (ii) Fig. \ref{fig:prompt_neg41}, \ref{fig:prompt_neg42}, \ref{fig:prompt_neg43}. We first leverage the bounding boxes and dense captions as the "visual" input. We provide the human instructions and responses from different models in Fig. \ref{fig:prompt_neg32} and Fig. \ref{fig:prompt_neg42}. Furthermore, we ask GPT4 to pretend as a smart teacher and score (0-10) the answers according to the image content and instructions. There are two criteria. \textit{(1) Accuracy: whether the response is accurate concerning the image content}. \textit{(2) Relevancy: whether the response directly follows the instruction}. After that, GPT4 is required to generate a score and a reason in order. Fig. \ref{fig:prompt_neg33}  and Fig. \ref{fig:prompt_neg43} show the full evaluation output from GAVIE.


% \section{LRV-Instruction with Pseudo Dense Captions}
% % \linjie{This could be a standalone section, should not be under assessment of GAVIE. The purpose is to show that our data generation pipeline can be extensible to arbitrary images with pseudo dense caption and bb annotation, for further scaling up. It is not for proving the robustness of GAVIE.....} \kevin{Agree. I have similar comments}
% % As \textit{GAVIE} takes the groundtruth dense caption annotations as inputs for evaluation, one may wonder whether \textit{GAVIE} can be applied to other datasets that do not have dense caption annotations.
% % To showcase the versatility of our proposed evaluation method, \textit{GAVIE}, we demonstrate its applicability to datasets that lack groundtruth dense caption annotations.
% To showcase the potential of scaling up LRV-Instruction to more images, especially the ones without dense caption annotations, we explore a pseudo labeling pipeline with 
% % To demonstrate the flexibility of the proposed \textit{GAVIE}, 
% % In the experiment, we employ
% a state-of-art model GRiT \cite{wu2022grit} to generate dense captions for the given images. % which are used as the pseudo groundtruths during evaluation. % the image content in the input prompt.
% These generated annotations can be directly used as the image content in the input prompt to GPT4.
% % the pseudo groundtruths during the evaluation process.
% As shown in Tab. \ref{tab:appendix_grit}, we compare the results obtained when using groundtruth dense captions from human annotation with those obtained when using % model-generated captions 
% the generated pseudo dense captions %as the image content 
% to evaluate the large multi-modal models (LMM) on our evaluation set. Remarkably, we observe that using the pseudo dense captions %model-generated image content 
% from GRiT yields similar results to using groundtruth dense captions. % as the image content.
% This serves as evidence of the extensibility of LRV-Instruction to images without dense caption annotations. 
% % the efficacy of utilizing \textit{GAVIE} for datasets that lack groundtruth dense caption annotations.
% % This demonstrates the effectiveness of using \textit{GAVIE} %model-generated content 
% % for datasets that lack groundtruth dense caption annotations. % This is because GRiT can generate high-quality bounding boxes and dense captions on average in the evaluation set, which can provide enough image information. This experiment demonstrates that \textit{GAVIE} can be potentially applied to other datasets without the groundtruth dense caption annotations.

% \begin{table}[h]
% \setlength\tabcolsep{3pt}
% \centering
% \small
% \begin{tabular}{lcccccc}
% \toprule[1.5pt]
% Image Content & Metric&Ours & MiniGPT4 & LLaVA & InstructBLIP & MMGPT\\
% \midrule
% GT Dense Captions &\textsc{Accuracy(GPT4)} & 6.58 & 4.14 & 4.36 & 5.93 & 0.91\\
% GT Dense Captions &\textsc{Relevancy(GPT4)} & 8.46 & 5.81 & 6.11 & 7.34 & 1.79\\
% \midrule
% Pseudo Dense Captions &\textsc{Accuracy(GPT4)} & 6.45 & 4.73 & 4.25 & 5.82 & 1.10\\
% Pseudo Dense Captions &\textsc{Relevancy(GPT4)} & 8.95 & 6.04 & 6.31 & 7.77 & 1.86\\
% \bottomrule
% \end{tabular}
% \vspace{0.06in}
% \caption{Comparison between using groundtruth (GT) and pseudo dense captions as the image content for LRV-Instruction data generation. The metric scores
% of \textsc{Accuracy(GPT4)} and \textsc{Relevancy(GPT4)} from our proposed GAVIE are in the scale of 0 to 10.
% }
% \label{tab:appendix_grit}
% \vspace{-0.2in}
% \end{table}


\subsection{\textit{GPT4-Assisted Visual Instruction Evaluation (GAVIE)} vs. Human Evaluation}
% \subsubsection{Replace Groundtruth Dense Captions with Model-Generated Dense Captions}

% \subsection{Human Evaluation}
In this section, we gain insights into the \textit{GAVIE} via human evaluation. Here, we randomly select 40 image-instruction instances from the evaluation set. The human evaluation is carried out by three experts specializing in the field of NLP. The questionnaire consists of 40 questions, which are randomly shuffled for each expert. The questionnaire takes about 20 minutes to complete on average. Each question includes an instruction, an image, and responses from 4 different LMMs. We provide instructions for experts as follows: 

\textit{"As for each question, there are an instruction, an image, and several answers. Suppose you are a smart teacher, please score the answers according to the two criteria. (1) Accuracy: whether the response is accurate concerning the image content. (2) Relevancy: whether the response directly follows the instruction without unrelated answers. There are four options for the scores (1) Very Poor, (2) Poor, (3) Good, (4) Excellent."}

\begin{table}[h]
\setlength\tabcolsep{3pt}
\centering
\small
\begin{tabular}{lccccc}
\toprule[1.5pt]
Evaluator& Ours & MiniGPT4 & LLaVA & InstructBLIP & MMGPT\\
\midrule
\textit{Expert1(1-4)}& \textcolor{red}{3.48} & \textcolor{blue}{2.61} & \textcolor{green}{2.87} & \textcolor{orange}{3.00} & \textcolor{magenta}{1.90}\\
\textit{Expert2(1-4)} &\textcolor{red}{3.58}& \textcolor{green}{2.23}& \textcolor{blue}{2.07}& \textcolor{orange}{2.48}& \textcolor{magenta}{1.05}\\
\textit{Expert3(1-4)} & \textcolor{red}{3.33} & \textcolor{blue}{2.58} & \textcolor{green}{2.92} & \textcolor{orange}{2.89} & \textcolor{magenta}{1.38}\\
\midrule
\textit{GAVIE-Accuracy (0-10)} & \textcolor{red}{6.58} & \textcolor{blue}{4.14} & \textcolor{green}{4.36} & \textcolor{orange}{5.93} & \textcolor{magenta}{0.91}\\
\textit{GAVIE-Relevancy (0-10)}  & \textcolor{red}{8.46} & \textcolor{blue}{5.81} & \textcolor{green}{6.11} & \textcolor{orange}{7.34} & \textcolor{magenta}{1.79}\\
\bottomrule
\end{tabular}
\vspace{0.06in}
\caption{GAVIE vs. Human Evaluation. GAVIE scores roughly align with the expert ratings. Numbers highlighted with \textcolor{red}{red}, \textcolor{orange}{orange}, \textcolor{green}{green}, \textcolor{blue}{blue}, and \textcolor{magenta}{magenta} indicate rank 1 to 5.
% \linjie{please include the GAVIE accuracy/relevancy score for all these models on the 40 examples, to show that GAVIE score roughly aligns with expert rating.}
}
\label{tab:appendix_human}
\vspace{-0.2in}
\end{table}

A questionnaire sample can be found in this \href{https://forms.office.com/r/uFSJTPCfH3}{link}. To evaluate the results quantitatively, we assign different scores for the options: \textit{Very Poor=1, Poor=2, Good=3, Excellent=4}. From Tab. \ref{tab:appendix_human}, all experts agree that the output from our model is the best, followed by InstructBLIP in second place, and MMGPT performs the worst. % is the second and Multimodal GPT performs worst, which is the same as the observation in Tab. \ref{tab:appendix_grit}. 
The observation is similar to that of \textit{GAVIE} evaluation results. % as shown in Tab. \ref{tab:appendix_grit}.
Although the ranking orders of MiniGPT4 and LLaVA from experts are not always the same as that of \textit{GAVIE}, the scores assigned to them are fairly close. One possible reason is that the answers from MiniGPT4 and LLaVA % are generally long, 
tend to be longer, % which is 
making them more challenging for humans to evaluate. % In general, the GPT4-Assisted evaluation approach can achieve the same level of results as human experts.
% Overall, 



\begin{figure*}[h]
    \centering
    % \vspace{-20mm}
      \centerline{\includegraphics[width=1\textwidth]{fuxiao-pos-example1.pdf}}
    \caption{The first example for generating positive instruction and answers (Part1). The image is not included in the text prompt. }
    \label{fig:prompt_pos11}
\end{figure*}

\begin{figure*}[h]
    \centering
      \includegraphics[width=1\textwidth]{fuxiao-pos-example1-result.pdf}
    \caption{The first example for generating positive instruction and answers (Part2).}
    \label{fig:prompt_pos12}
\end{figure*}

\begin{figure*}[h]
    \centering
      \includegraphics[width=1\textwidth]{fuxiao-pos-example1-result2.pdf}
    \caption{The first example for generating positive instruction and answers (Part3).}
    \label{fig:prompt_pos13}
\end{figure*}

\begin{figure*}[h]
    \centering
      \includegraphics[width=1\textwidth]{fuxiao-pos-example2.pdf}
    \caption{The second example for generating positive instruction and answers (Part1). The image is not included in the text prompt.}
    \label{fig:prompt_pos21}
\end{figure*}

\begin{figure*}[h]
    \centering
      \includegraphics[width=1\textwidth]{fuxiao-pos-example2-result.pdf}
    \caption{The second example for generating positive instruction and answers (Part2).}
    \label{fig:prompt_pos22}
\end{figure*}

\begin{figure*}[h]
    \centering
      \includegraphics[width=1\textwidth]{fuxiao-pos-example2-result2.pdf}
    \caption{The second example for generating positive instruction and answers (Part3).}
    \label{fig:prompt_pos23}
\end{figure*}

\begin{figure*}[h]
    \centering
      \includegraphics[width=1\textwidth]{fuxiao-neg1-example1.pdf}
    \caption{The example for generating negative instruction and answers by nonexistent element manipulation (Part1). The image is not included in the text prompt.}
    \label{fig:prompt_neg11}
\end{figure*}

\begin{figure*}[h]
    \centering
      \includegraphics[width=1\textwidth]{fuxiao-neg1-example-result.pdf}
    \caption{The example for generating negative instruction and answers by nonexistent element manipulation (Part2).}
    \label{fig:prompt_neg12}
\end{figure*}

\begin{figure*}[h]
    \centering
      \includegraphics[width=1\textwidth]{fuxiao-neg2-example1.pdf}
    \caption{The example for generating negative instruction and answers by existent element manipulation (Part1). The image is not included in the text prompt.}
    \label{fig:prompt_neg21}
\end{figure*}

\begin{figure*}[h]
    \centering
      \includegraphics[width=1\textwidth]{fuxiao-neg2-example-result2.pdf}
    \caption{The first example for generating negative instruction and answers by existent element manipulation (Part2).}
    \label{fig:prompt_neg22}
\end{figure*}

\begin{figure*}[h]
    \centering
      \includegraphics[width=1\textwidth]{fuxiao-eval1-example1.pdf}
    \caption{The first example for the GPT4-Assisted Visual Instruction Evaluation (Part1). The image is not included in the text prompt.}
    \label{fig:prompt_neg31}
\end{figure*}

\begin{figure*}[h]
    \centering
      \includegraphics[width=1\textwidth]{fuxiao-eval1-example1-xia.pdf}
    \caption{The first example for the GPT4-Assisted Visual Instruction Evaluation (Part2).}
    \label{fig:prompt_neg32}
\end{figure*}

\begin{figure*}[h]
    \centering
      \includegraphics[width=1\textwidth]{fuxiao-eval1-example1-result.pdf}
    \caption{The first example for the GPT4-Assisted Visual Instruction Evaluation (Part3).}
    \label{fig:prompt_neg33}
\end{figure*}

\begin{figure*}[h]
    \centering
      \includegraphics[width=1\textwidth]{fuxiao-eval2-example1.pdf}
    \caption{The second example for the GPT4-Assisted Visual Instruction Evaluation (Part1). The image is not included in the text prompt.}
    \label{fig:prompt_neg41}
\end{figure*}

\begin{figure*}[h]
    \centering
      \includegraphics[width=1\textwidth]{fuxiao-eval2-example2.pdf}
    \caption{The second example for the GPT4-Assisted Visual Instruction Evaluation (Part2).}
    \label{fig:prompt_neg42}
\end{figure*}

\begin{figure*}[h]
    \centering
      \includegraphics[width=1\textwidth]{fuxiao-eval2-example-result.pdf}
    \caption{The second example for the GPT4-Assisted Visual Instruction Evaluation (Part3).}
    \label{fig:prompt_neg43}
\end{figure*}


\begin{figure}
     \centering
     \begin{subfigure}[b]{\textwidth}
         \centering
         \includegraphics[width=0.8\textwidth]{neg11.pdf}
         \caption{All LMMs except ours produce inconsistent descriptions with the images. MiniGPT4 also has repetitive sentences.}
         \label{fig:NEG11}
     \end{subfigure}
     \par\bigskip
     \begin{subfigure}[b]{\textwidth}
         \centering
         \includegraphics[width=0.8\textwidth]{neg12.pdf}
         \caption{MiniGPT4, LLaVA, and InstructBLIP respond with long descriptive answers about the "painting", which does not exist in the image. Multimodal-GPT produces a question about the "painting," but it does not address the instruction. }
         \label{fig:NEG13}
     \end{subfigure}
     \par\bigskip
     \begin{subfigure}[b]{\textwidth}
         \centering
         \includegraphics[width=0.8\textwidth]{neg13.pdf}
         \caption{MiniGPT4, LLaVA, and Multimodal-GPT respond with a long descriptive answer about the "pink flowers", which do not exist in the image. Although InstructBLIP's answer is brief, it's still false.}
         \label{fig:NEG12}
     \end{subfigure}
        \caption{Model predictions on negative instruction examples of Neg1:\textit{"Nonexistent Element Manipulation"}. \textcolor{red}{RED} represents the nonexistent elements in the images.}
        \label{fig:demo_neg1}
\end{figure}
\begin{figure}
     \centering
     \begin{subfigure}[b]{\textwidth}
         \centering
         \includegraphics[width=0.9\textwidth]{neg21.pdf}
         \caption{All LMMs except ours produce inconsistent descriptions with the images. Our model can point out that the loose strands do not form a butterfly pattern. }
         \label{fig:NEG21}
     \end{subfigure}
     \par\bigskip
     \begin{subfigure}[b]{\textwidth}
         \centering
         \includegraphics[width=0.9\textwidth]{neg22.pdf}
         \caption{All LMMs except ours produce false attributes of windows in the images. Our model can point out that the windows are white instead of red. }
         \label{fig:NEG22}
     \end{subfigure}
     \par\bigskip
     \begin{subfigure}[b]{\textwidth}
         \centering
         \includegraphics[width=0.9\textwidth]{neg23.pdf}
         \caption{All LMMs except ours produce inconsistent descriptions with the images. Our model can point out that the woman with green hair doesn't have a black backpack on her shoulder. }
         \label{fig:NEG23}
     \end{subfigure}
        \caption{Model predictions on negative instruction examples of Neg2:\textit{"Existent Element Manipulation"}. \textcolor{red}{RED} represents the wrong attributes of existent objects in the images.}
        \label{fig:demo_neg2}
        
\end{figure}

\begin{figure}
     \centering
     \begin{subfigure}[b]{\textwidth}
         \centering
         \includegraphics[width=0.9\textwidth]{pos1.pdf}
         \caption{Object Detection Task. MiniGPT4, LLaVa, and MMGPT can't locate the objects correctly.}
         \label{fig:pos1}
     \end{subfigure}
     \par\bigskip
     \begin{subfigure}[b]{\textwidth}
         \centering
         \includegraphics[width=0.9\textwidth]{pos2.pdf}
         \caption{Multi-choice VQA Task. InstructBLIP can't follow the instruction well by generating anything. LLaVA's answer seems reasonable, but it incorrectly detects the clouds in the image. }
         \label{fig:pos2}
     \end{subfigure}
     \par\bigskip
     \begin{subfigure}[b]{\textwidth}
         \centering
         \includegraphics[width=0.9\textwidth]{pos3.pdf}
         \caption{Referential Expression Grounding Task. Other LMMs can't directly address the instruction by giving a general answer, "kitchen," while our model can clearly predict the "floor". }
         \label{fig:pos3}
     \end{subfigure}
        \caption{Positive instruction demos with different tasks and predictions from different models.}
        \label{fig:demo_pos1}
\end{figure}


\begin{figure}
     \centering
     \begin{subfigure}[b]{\textwidth}
         \centering
         \includegraphics[width=0.9\textwidth]{pos4.pdf}
         \caption{Image Anomaly Detection Task. Our model and LLaVA point out there is a normal image without unusual elements, while MiniGPT4, InstructBLIP, and MMGPT describe the image instead of addressing the instruction.}
         \label{fig:demo2_pos1}
     \end{subfigure}
     \par\bigskip
     \begin{subfigure}[b]{\textwidth}
         \centering
         \includegraphics[width=0.9\textwidth]{pos5.pdf}
         \caption{Visual Entailment Task. }
         \label{fig:demo2_pos2}
     \end{subfigure}
     \par\bigskip
     \begin{subfigure}[b]{\textwidth}
         \centering
         \includegraphics[width=0.9\textwidth]{pos6.pdf}
         \caption{Object Interaction Analysis Task. All LMMs except ours describe the wrong location of the book. LLaVA generates long text with unrelated information to address the instruction.}
         \label{fig:demo2_pos3}
     \end{subfigure}
        \caption{Positive instruction demos with different tasks and predictions from different models. }
        \label{fig:demo_pos2}
\end{figure}


\begin{figure}
     \centering
     \begin{subfigure}[b]{\textwidth}
         \centering
         \includegraphics[width=0.9\textwidth]{pos7.pdf}
         \caption{Dense Description Task. All LMMs except ours incorrectly recognize the color of the bridle. InstructBLIP also fails to answer the "state of the horse's bridle" in the instruction.}
         \label{fig:demo3_pos1}
     \end{subfigure}
     \par\bigskip
     \begin{subfigure}[b]{\textwidth}
         \centering
         \includegraphics[width=0.9\textwidth]{pos8.pdf}
         \caption{Image Caption Task. LLaVA and InstructBLP fail to follow the instruction for generating one caption. Instead, they generate long text with several sentences.}
         \label{fig:demo3_pos2}
     \end{subfigure}
     \par\bigskip
     \begin{subfigure}[b]{\textwidth}
         \centering
         \includegraphics[width=0.9\textwidth]{pos9.pdf}
         \caption{Activity Recognition Task.}
         \label{fig:demo3_pos3}
     \end{subfigure}
        \caption{Positive instruction demos with different tasks and predictions from different models. }
        \label{fig:demo_pos3}
\end{figure}
\label{sec:appendix}




\end{document}



